\section{Global invariant descent on the certified finite atlas}
\label{sec:block13-global-invariant}

Global synthesis atlas performs the assembly step requested after the native stratum invariants.  The preceding sections do not justify one unrestricted equivalence among all eight broad strata.  They do, however, supply several finite connected systems of normalized charts for which the semantic object, overlap categories, bridge directions, information loss, coherence, and local completeness are explicit.  This section freezes those systems into one authoritative certified atlas, constructs the invariant comparison maps, verifies their higher coherence, and proves the strongest global completeness theorem supported by the assembled results.

The global object in this section is deliberately a finite disjoint union of connected semantic components.  Each component contains nontrivial reversible bridges across different kinds of presentation.  The component tag is retained as part of the invariant.  No bridge is invented between two components merely because their formulas use similar words or symbols.  In particular, the functional-Mahler coefficient component and the continuous Newton--Mahler interpolation component remain distinct even though both belong to the broad local stratum.

\subsection{The authoritative bridge matrix}
\label{subsec:block13-bridge-matrix}

Write
\[
 \mathsf E,
 \quad
 \mathsf{FF},
 \quad
 \mathsf D,
 \quad
 \mathsf P,
 \quad
 \varnothing
\]
for an equivalence on a named normalized overlap, a fully faithful inclusion, an information-losing directional functor, a partial or hypothesis-dependent bridge, and the absence of a certified broad bridge, respectively.  The matrix below records only relations already established in the preceding sections or in the audited source concordance.  A symbol carrying a subscript refers to the named overlap rather than to the whole broad stratum.

\begin{landscape}
\begin{table}[p]
\centering
\caption{Authoritative broad-stratum bridge matrix for the certified global synthesis atlas. Every decorated equivalence is restricted to the named normalized overlap; the undecorated broad categories are not thereby identified.}
\label{tab:block13-bridge-matrix}
\scriptsize
\setlength{\tabcolsep}{3.7pt}
\begin{tabular}{@{}lcccccccc@{}}
\toprule
Source $\backslash$ target & LIN & HOL & DIG & ALG & TREE & PROD & LOC & GAL\\
\midrule
LIN  & $\mathsf E$ & $\mathsf{FF}$ & $\mathsf E_{\rm per}$ & -- & -- & -- & $\mathsf E_{\rm per}$ & --\\
HOL  & $\mathsf P$ & $\mathsf E$ & -- & $\mathsf E_{\rm Fuss}/\mathsf P$ & $\mathsf E_{\rm Fuss}/\mathsf P$ & -- & -- & --\\
DIG  & $\mathsf E_{\rm per}$ & $\mathsf P$ & $\mathsf E$ & $\mathsf P$ & -- & -- & $\mathsf E_{\rm rm}$ & --\\
ALG  & -- & $\mathsf E_{\rm Fuss}/\mathsf D$ & $\mathsf P$ & $\mathsf E$ & $\mathsf E_{\rm Fuss}/\mathsf P$ & -- & -- & --\\
TREE & -- & $\mathsf E_{\rm Fuss}/\mathsf D$ & -- & $\mathsf E_{\rm Fuss}/\mathsf D$ & $\mathsf E$ & -- & -- & --\\
PROD & $\mathsf D$ & -- & -- & -- & -- & $\mathsf E$ & $\mathsf P$ & --\\
LOC  & $\mathsf E_{\rm per}$ & -- & $\mathsf E_{\rm rm}$ & -- & -- & $\mathsf P$ & $\mathsf E$ & --\\
GAL  & -- & -- & -- & -- & -- & -- & -- & $\mathsf E$\\
\bottomrule
\end{tabular}
\vspace{5pt}
\parbox{0.96\linewidth}{\scriptsize
\textbf{Legend.} $\mathsf E$: equivalence; $\mathsf{FF}$: fully faithful inclusion; $\mathsf D$: information-losing directional functor; $\mathsf P$: partial, constructor-dependent, or certificate-dependent bridge; --: no certified bridge. The subscripts $\mathrm{per}$, $\mathrm{rm}$, and $\mathrm{Fuss}$ denote the unit-periodic, rational functional-Mahler, and marked weighted Fuss overlaps. A cell of the form $\mathsf E_{\mathrm{Fuss}}/\mathsf D$ means equivalence on the marked Fuss slice and directional loss on the ambient bridge. Functional-Mahler coefficient equations and Newton--Mahler interpolation remain distinct despite sharing the LOC row.}
\end{table}
\end{landscape}


The matrix has three immediate consequences.

\begin{enumerate}
\item The reversible core is strictly smaller than the directed transformation category.
\item A bridge may become an equivalence after restricting to a normalized overlap without becoming invertible on the ambient stratum.
\item Empty cells are mathematical data.  They prevent later arguments from silently transporting an invariant through an uncertified route.
\end{enumerate}

The entries \(E_{\mathrm{per}}\), \(E_{\mathrm{rm}}\), and \(E_{\mathrm{mark}}\) refer, respectively, to the unit-periodic overlap, the rational functional-Mahler overlap, and the marked polynomial fixed-point overlap.  In particular, the functional-Mahler and Newton--Mahler languages remain separate outside their explicitly declared interfaces, and faithful corner realization remains separate from neutral Tannakian reconstruction.

The full object-level bridge ledger, including every normalized source, target, overlap, invariant comparison, and loss or gate, is retained as a machine-generated source artifact.  The printed matrix and the component table below contain the proof-relevant summary; a relation absent from the ledger is not used in the descent nerve.

\subsection{Certified connected components}
\label{subsec:block13-components}

\begin{definition}[Certified Global synthesis atlas atlas]
\label{def:block13-certified-atlas}
Fix a prime \(p\), a characteristic-zero field \(K\), and a finite Galois extension \(L/K\) for the Galois component.  The certified atlas \(\mathfrak A_{13}\) is the finite coproduct of the following five connected profiled atlases.

\begin{enumerate}[label=\textup{(C\arabic*)}]
\item \(\mathfrak A_{\mathrm{DRLM},p}\), the fixed-prime digital--regular component with rational functional-Mahler and periodic linear subcharts;
\item \(\mathfrak A_{\mathrm{Loc},p}\), the continuous Newton--Mahler and complete clopen-branch component;
\item \(\mathfrak A_{\mathrm{Fuss}}\), the marked Fuss--Catalan nonlinear--algebraic--holonomic component;
\item \(\mathfrak A_{\mathrm{EM}}\), the fixed-Euler-basis and prime-uniform multiplicative component;
\item \(\mathfrak A_{\mathrm{Gal},L/K}\), the finite semisimple Galois component for the fixed extension \(L/K\).
\end{enumerate}

The reversible arrows are exactly the normalized equivalences printed below.  Directional constructions, reductions, forgetful maps, and obstructed lifts are attached as satellites but do not enter the maximal subgroupoid used for invariant descent.
\end{definition}

\begin{table}[htbp]
\centering
\caption{Certified connected components of the Global synthesis atlas atlas.}
\label{tab:block13-components}
\tiny
\begin{tabularx}{\textwidth}{@{}L{0.105\textwidth}L{0.255\textwidth}L{0.285\textwidth}Y@{}}
\toprule
Component & Semantic domain and normalized charts & Complete coordinate and reversible overlap & Directional satellites or exact boundary\\
\midrule
DRLM & fixed-prime regular sequences; canonical section module, strict shortlex record, and minimal DFAO; rational-mask/Cartier and periodic linear subcharts & strict shortlex pointed digit-linear record, refined on the rational-mask and unit-periodic full replete subcategories & integral reductions and Newton--Mahler interpolation are directional; many automatic sequences have no rational mask\\
LOC & continuous $\mathbb Z_p$-valued functions; Newton--Mahler coefficients, complete clopen branch packets, and affine branch coordinates & canonical Mahler coefficient sequence; complete branch analysis and synthesis at each declared depth & functional-Mahler equations require a separately declared coefficient lift and interpolation constructor\\
Fuss & marked $U=z+\alpha U^q$ tree system, selected algebraic branch, and compressed P-recursive coefficient chart & parameter pair $(q,\alpha)$ with its support and singularity-orbit package & forgetting the branch mark or generator grammar is information-losing\\
EM & fixed-basis Euler exponents, divisor ghosts, integral unit series, and prime-uniform multiplicative local factors & Euler exponent sequence on the fixed-basis, prime-uniform essential image & arbitrary product grammars and nonuniform prime-local fields remain outside the reversible overlap\\
Gal & subgroup, fixed field, and faithful Reynolds projector for one fixed finite Galois extension & actual subgroup $H\leq G$ with functorial core, normal-closure, and double-coset data & projective Brauer obstruction, modular or wild failure, and the tensor-fiber gate remain separately typed\\
\bottomrule
\end{tabularx}
\end{table}


\begin{figure}[htbp]
\centering
\includegraphics[width=.99\textwidth]{global_certified_atlas.pdf}
\caption{The certified finite atlas.  Each box is a nontrivial connected system of normalized charts.  Dashed external arrows are directional, partial, or obstructed satellites and are not promoted to equivalences.}
\label{fig:block13-certified-atlas}
\end{figure}

The semantic category of the total atlas is
\[
 \mathcal O_{13}
 =
 \mathcal O_{\mathrm{DRLM}}
 \sqcup
 \mathcal O_{\mathrm{Loc}}
 \sqcup
 \mathcal O_{\mathrm{Fuss}}
 \sqcup
 \mathcal O_{\mathrm{EM}}
 \sqcup
 \mathcal O_{\mathrm{Gal}}.
\]
The coproduct symbol is literal.  It does not assert that a finite-field regular sequence, a continuous \(p\)-adic function, a characteristic-zero marked tree branch, an integral Euler template, and a finite Galois subgroup are objects of one untyped category.  The tag records which semantic doctrine is in force.

\subsection{Normalizer-induced higher coherence}
\label{subsec:block13-normalizer-coherence}

Every connected component in \cref{def:block13-certified-atlas} has a canonical semantic normalizer.  For a chart \(\mathcal C_{\lambda,a}\) in component \(\lambda\), write
\[
 N_{\lambda,a}:\mathcal C_{\lambda,a}\longrightarrow\mathcal O_\lambda
\]
for the semantic normalization functor and choose its canonical inverse chart
\[
 Q_{\lambda,a}:\mathcal O_\lambda\longrightarrow\mathcal C_{\lambda,a}.
\]
The certified bridge from chart \(a\) to chart \(b\) is
\[
 \Phi_{ab}=Q_{\lambda,b}N_{\lambda,a}.
\]

\begin{theorem}[Normalizer-induced coherence]
\label{thm:block13-normalizer-coherence}
Assume that every normalized fiber of \(N_{\lambda,a}\) is contractible.  Then the comparison isomorphism
\[
 \eta_{abc}:
 \Phi_{bc}\Phi_{ab}
 \Longrightarrow
 \Phi_{ac}
\]
is unique.  Every triple-overlap diagram commutes, every quadruple-overlap coherence equation holds, and the composite around every finite bridge cycle is the identity transformation of the normalized chart.
\end{theorem}

\begin{proof}
Both \(\Phi_{bc}\Phi_{ab}(X)\) and \(\Phi_{ac}(X)\) are normalized presentations of the same semantic object \(N_{\lambda,a}(X)\).  Contractibility of the normalized fiber gives a unique isomorphism between them.  Any two composites occurring in a coherence diagram are isomorphisms between the same two normalized presentations, so uniqueness makes them equal.  The cycle statement follows by applying the same argument to the starting chart and itself.
\end{proof}

\begin{figure}[htbp]
\centering
\includegraphics[width=.93\textwidth]{normalizer_induced_coherence.pdf}
\caption{Every reversible chart in one certified component normalizes to the same semantic object.  The unique normalized comparison forces all finite bridge cycles to commute.}
\label{fig:block13-normalizer-coherence}
\end{figure}

This theorem is not a substitute for constructing the normalizer.  It applies here because the preceding sections have already supplied unique reduced masks, reachable-observable Cartier modules, minimal automata, canonical section modules, marked fixed-point systems, selected algebraic branches, fixed Euler exponents, prime-power local factors, fixed subgroups, fixed fields, and faithful Reynolds projectors on the declared domains.

\subsection{Invariant descent theorem}
\label{subsec:block13-invariant-descent}

For every chart \(\mathcal C_{\lambda,a}\), let
\[
 I_{\lambda,a}:\mathcal C_{\lambda,a}\longrightarrow\mathcal V_{\lambda,a}
\]
be its native normalized invariant.  A bridge comparison is a natural isomorphism
\[
 \kappa_{ab}:
 J_{ab}I_{\lambda,a}
 \Longrightarrow
 I_{\lambda,b}\Phi_{ab},
\]
where \(J_{ab}\) transports invariant coordinates between their chart-specific target categories.

\begin{definition}[Coherent invariant family]
The family \((I_{\lambda,a},\kappa_{ab})\) is coherent when the two induced comparison maps from \(a\) to \(c\) agree on every triple overlap and the corresponding quadruple-overlap equation holds.
\end{definition}

\begin{theorem}[Certified finite invariant descent]
\label{thm:block13-certified-invariant-descent}
Let \(\mathfrak A_{13}\) be the atlas of \cref{def:block13-certified-atlas}.  Suppose that, in every connected component:
\begin{enumerate}
\item the local invariant is complete on every normalized chart;
\item the bridge comparisons are coherent;
\item invariant-compatible local isomorphisms satisfy effective descent.
\end{enumerate}
Then the local invariants induce a functor
\[
 \mathcal I_{13}:
 \Desc(\mathfrak A_{13})
 \longrightarrow
 \Desc(\mathfrak V_{13})
\]
that reflects isomorphisms.  If each semantic component is discrete, equality of descended invariants is equivalent to equality of the represented semantic objects.
\end{theorem}

\begin{proof}
The local comparison maps form a pseudonatural invariant family, so \cref{thm:atlas-assembly} gives the descended functor.  Equality or isomorphism of the descended invariant gives compatible local invariant matches.  Local completeness supplies local object isomorphisms.  Effective isomorphism descent glues those local isomorphisms to an isomorphism in the descent category.  On a discrete semantic component, every isomorphism is equality.
\end{proof}

\begin{figure}[htbp]
\centering
\includegraphics[width=.95\textwidth]{global_invariant_descent.pdf}
\caption{The global invariant theorem has four separate requirements.  Local completeness alone is insufficient without comparison maps, higher coherence, and effective descent of the resulting isomorphisms.}
\label{fig:block13-invariant-descent}
\end{figure}

\subsection{Digital--regular--linear--functional-Mahler component}
\label{subsec:block13-drlm}

Fix the prime \(p\) and let
\[
 \mathcal O_{\mathrm{DRLM}}
 =\mathbf{Reg}_p(\mathbb F_p)
\]
be the discrete semantic category of \(p\)-regular sequences over \(\mathbb F_p\).  By the finite-field fixed-radix theorem, every object is also \(p\)-automatic.  Its globally defined normalized charts are the canonical section module, the strict shortlex pointed digit-linear record, and the minimal least-significant-digit-first \(p\)-DFAO.

Two full replete subcategories carry additional charts.
\begin{enumerate}
\item On the rational-mask locus \(\RMah_p^{\times}\), the reduced functional-Mahler mask and the minimal reachable-observable Cartier module are equivalent to the digital charts.
\item On the unit purely periodic locus, the minimal recurrence, reduced rational OGF, minimal finite linear state, primitive period word, and rational-mask chart are all equivalent.
\end{enumerate}

For \(a\in\mathcal O_{\mathrm{DRLM}}\), define
\[
 \mathcal J_{\mathrm{DRLM},p}(a)
 =
 \left(
 \mathsf{SL}_{p,\mathbb F_p}(a),
 \mathcal P_{\mathrm{rm}}(a),
 \mathcal P_{\mathrm{per}}(a)
 \right),
\]
where \(\mathsf{SL}_{p,\mathbb F_p}(a)\) is the strict shortlex record, \(\mathcal P_{\mathrm{rm}}(a)\) is empty off the rational-mask locus and otherwise contains the reduced mask, Cartier module, collapse obstruction, and minimal automaton coordinates, and \(\mathcal P_{\mathrm{per}}(a)\) is empty off the purely periodic locus and otherwise contains the primitive period word, exact period, minimal recurrence, reduced OGF, and periodic packet profile.

\begin{theorem}[Digital--regular--linear--Mahler comparison theorem]
\label{thm:block13-drlm-comparison}
The invariant \(\mathcal J_{\mathrm{DRLM},p}\) is preserved by every reversible bridge in \(\mathfrak A_{\mathrm{DRLM},p}\).  Its comparison cells satisfy all higher coherence equations.  The strict shortlex coordinate reconstructs the represented regular sequence, so the descended invariant is complete on \(\mathbf{Reg}_p(\mathbb F_p)\).  The mask and periodic packages are compatible exact refinements on their respective subloci.
\end{theorem}

\begin{proof}
The canonical section module, strict shortlex record, and minimal orbit automaton normalize to the same coefficient sequence by the Fixed-radix regular atlas effective-descent theorem.  On the rational-mask locus, fixed-denominator Cartier intertwining identifies the reachable numerator span with the canonical section module and the reachable numerator set with the minimal automaton states.  On the periodic locus, the recurrence, rational, state, automaton, and mask charts normalize to the same primitive coefficient word.  Every comparison therefore factors through the same semantic normalizer, and \cref{thm:block13-normalizer-coherence} applies.  Completeness follows from the strict shortlex reconstruction formula.
\end{proof}

\begin{table}[htbp]
\centering
\caption{Representative cross-stratum invariant comparisons.}
\label{tab:block13-examples}
\scriptsize
\begin{tabularx}{\textwidth}{@{}L{0.15\textwidth}L{0.12\textwidth}L{0.13\textwidth}L{0.13\textwidth}L{0.13\textwidth}Y@{}}
\toprule
Example & Period & Recurrence order & Regular rank & DFAO states & Mask compiler degree\\
\midrule
$\mathbb {F}_2$/ word 10 & 2 & 2 & 2 & 3 & 2\\
$\mathbb {F}_2$/ word 110 & 3 & 2 & 2 & 3 & 2\\
$\mathbb {F}_3$/ word 12 & 2 & 1 & 1 & 2 & 1\\
\bottomrule
\end{tabularx}
\end{table}


The table gives representative objects for which recurrence order, regular rank, automaton size, and mask-compiler dimension are different coordinates of the same normalized semantic sequence.  Equality of one displayed scalar alone is not a complete invariant; the complete records are the chart objects described above.

The inclusions of the rational-mask and periodic loci are not essentially surjective.  Over a finite field every regular sequence is automatic, but many automatic sequences have no rational canonical mask, and many rational-mask sequences are not periodic.  The optional fields in \(\mathcal J_{\mathrm{DRLM},p}\) record these exact domain boundaries rather than forcing one chart to exist everywhere.

\subsection{Continuous Newton--Mahler component}
\label{subsec:block13-loc}

Let
\[
 \mathcal O_{\mathrm{Loc}}
 =C(\mathbb Z_p,\mathbb Z_p)
\]
with equality morphisms.  Its canonical coordinate is the Mahler sequence
\[
 \mathcal M_p(f)
 =
 (c_r(f))_{r\ge0},
 \qquad
 c_r(f)=\Delta^r f(0),
 \qquad
 c_r(f)\longrightarrow0.
\]
At every declared depth \(p^a\), the complete clopen branch packet
\[
 \bigl(f(p^aq+s)\bigr)_{0\le s<p^a}
\]
and the corresponding branchwise Mahler coefficients are reversible by explicit piecewise synthesis.  Affine branch kernels and mixed-radix refinements are comparison charts of the same semantic function.

\begin{theorem}[Continuous local invariant comparison]
\label{thm:block13-local-comparison}
The Mahler coefficient sequence
\[
 \mathcal J_{\mathrm{Loc},p}(f)
 =
 (c_r(f))_{r\ge0}
\]
is a complete invariant of \(\mathfrak A_{\mathrm{Loc},p}\).  Complete branch packets, branchwise Mahler coordinates, affine coefficient kernels, and nested refinements admit coherent comparison maps to \(\mathcal J_{\mathrm{Loc},p}\).
\end{theorem}

\begin{proof}
Mahler analysis and synthesis are mutually inverse on \(C(\mathbb Z_p,\mathbb Z_p)\).  Complete clopen analysis and piecewise synthesis are mutually inverse at each depth, and the nested branch law is strict.  Conjugating those maps by Mahler analysis gives the affine coefficient comparisons, whose mixed-radix composition was proved in Local Mahler-polyphase atlas.  All normalized charts therefore return the same function and the same Mahler sequence, so \cref{thm:block13-normalizer-coherence} gives higher coherence.
\end{proof}

Functional-Mahler equations for coefficient series are not included in this component.  A finite-state or functional-Mahler object enters the Newton--Mahler component only through a separately declared lift and interpolation constructor carrying its own precision and information-loss certificate.

\subsection{A genuine nonlinear--algebraic--holonomic overlap}
\label{subsec:block13-fuss}

Let \(K\) be a characteristic-zero field.  For an integer \(q\geq2\) and \(\alpha\in K^\times\), let
\[
 U_{q,\alpha}(z)\in zK[[z]]
\]
be the unique selected branch satisfying
\begin{equation}
 U=z+\alpha U^q.
 \label{eq:block13-fuss-tree}
\end{equation}
The derivative of
\[
 P_{q,\alpha}(z,Y)=Y-z-\alpha Y^q
\]
with respect to \(Y\) is one at \((0,0)\), so the branch is simple and selected by \(U(0)=0\).

Put
\[
 x=z^{q-1},
 \qquad
 U_{q,\alpha}(z)=zC_{q,\alpha}(x).
\]
Then
\begin{equation}
 C=1+\alpha xC^q.
 \label{eq:block13-fuss-compressed}
\end{equation}
Lagrange inversion gives
\begin{equation}
 [x^n]C_{q,\alpha}(x)
 =
 c_n(q,\alpha)
 =
 \frac{\alpha^n}{(q-1)n+1}
 \binom{qn}{n}.
 \label{eq:block13-fuss-coefficients}
\end{equation}
Define
\[
 A_q(n)
 =(n+1)\prod_{r=2}^{q}\bigl((q-1)n+r\bigr),
\]
\[
 B_q(n)
 =\prod_{r=1}^{q}(qn+r).
\]
Then the coefficient sequence satisfies the first-order polynomial recurrence
\begin{equation}
 \boxed{
 A_q(n)c_{n+1}
 =
 \alpha B_q(n)c_n.}
 \label{eq:block13-fuss-recurrence}
\end{equation}
After dividing the polynomial gcd of \(A_q\) and \(B_q\), this is the normalized first-order tail recurrence with initial value \(c_0=1\).

\begin{theorem}[Fuss--Catalan three-stratum equivalence]
\label{thm:block13-fuss-equivalence}
The following normalized marked presentations determine one another exactly:
\begin{enumerate}
\item the polynomial tree system \(U=z+\alpha U^q\);
\item the selected algebraic branch \((P_{q,\alpha},U(0)=0)\);
\item the compressed P-recursive coefficient chart \((\cref{eq:block13-fuss-recurrence},c_0=1)\) together with the support embedding \(U(z)=zC(z^{q-1})\).
\end{enumerate}
Every normalized fiber is contractible.  All triangle comparisons return the same parameter pair \((q,\alpha)\).
\end{theorem}

\begin{proof}
The marked tree equation is the selected algebraic equation.  Simplicity at the origin gives its unique formal branch.  Substitution \(U=zC(z^{q-1})\) gives \cref{eq:block13-fuss-compressed}.  Lagrange inversion gives \cref{eq:block13-fuss-coefficients}, and taking the ratio of consecutive coefficients gives \cref{eq:block13-fuss-recurrence}.  Conversely, the recurrence and initial value reconstruct \(C\), the support embedding reconstructs \(U\), and the reconstructed series satisfies the selected algebraic equation.  A normalized automorphism fixes \(q\), \(\alpha\), and the unique branch, hence is the identity.
\end{proof}

\begin{figure}[htbp]
\centering
\includegraphics[width=.87\textwidth]{fuss_catalan_three_stratum.pdf}
\caption{The marked Fuss--Catalan family gives a genuine overlap among the nonlinear tree, selected algebraic, and holonomic recurrence strata.  Every path returns the same marked pair \((q,\alpha)\).}
\label{fig:block13-fuss-triangle}
\end{figure}

The tree arity period is \(q-1\).  For packet modulus \(N\), set
\[
 g_N=\gcd(N,q-1).
\]
The nonzero coefficient addresses of \(U_{q,\alpha}\) are exactly
\[
 j\equiv1\pmod{g_N},
\]
and the complete dominant singularity orbit has size
\[
 \frac{q-1}{g_N}
\]
under the positive analytic hypotheses of \cref{thm:block11-qary-phase}.  These packet signatures agree in the tree and algebraic charts because both are computed from the same selected branch.

Define
\[
 \mathcal J_{\mathrm{Fuss}}(q,\alpha)
 =
 \left(
 q,
 \alpha,
 P_{q,\alpha},
 \{c_n(q,\alpha)\}_{n\geq0},
 \{g_N\}_{N\geq1},
 \left\{\frac{q-1}{g_N}\right\}_{N\geq1}
 \right).
\]

\begin{corollary}[Complete cross-stratum invariant on the marked Fuss family]
\label{cor:block13-fuss-invariant}
The nonlinear, algebraic, holonomic, packet-support, and dominant-orbit invariants admit coherent comparison maps to \(\mathcal J_{\mathrm{Fuss}}\).  The package is complete on the declared marked family.
\end{corollary}

\begin{proof}
Coherence follows from \cref{thm:block13-fuss-equivalence,thm:block13-normalizer-coherence}.  The first two entries reconstruct the marked fixed-point equation and therefore its unique selected branch.
\end{proof}

\begin{table}[htbp]
\centering
\caption{Named weighted q-ary examples in the cross-stratum atlas.}
\label{tab:block13-qary-examples}
\scriptsize
\begin{tabularx}{\textwidth}{@{}L{0.10\textwidth}L{0.11\textwidth}L{0.12\textwidth}L{0.25\textwidth}L{0.17\textwidth}Y@{}}
\toprule
$q$ & $\lambda$ & Support period & First supported coefficients & Active addresses at $N=6$ & Packet-orbit size\\
\midrule
2 & $1$ & 1 & $(1, 1, 2, 5,\ldots)$ & $\{0,1,2,3,4,5\}$ & 1\\
3 & $2$ & 2 & $(1, 2, 12, 96,\ldots)$ & $\{1,3,5\}$ & 1\\
4 & $-1$ & 3 & $(1, -1, 4, -22,\ldots)$ & $\{1,4\}$ & 1\\
5 & $1/2$ & 4 & $(1, 1/2, 5/4, 35/8,\ldots)$ & $\{1,3,5\}$ & 2\\
\bottomrule
\end{tabularx}
\end{table}


The examples show that the finite invariant retains the weight parameter as well as the arity.  Support period alone recovers only \(q-1\); the coefficient \([z^q]U=\alpha\) supplies the second coordinate.

This is not a classification of arbitrary algebraic branches, arbitrary P-recursive sequences, or unmarked nonlinear fixed points.  It is the first explicit complete invariant on a connected overlap of those three strata.

\subsection{Euler--multiplicative component}
\label{subsec:block13-em}

Let
\[
 A_b(z)=\prod_{m\geq1}(1-z^m)^{-b_m}
\]
be an integral unit series in the fixed Euler basis, and let
\[
 c_n=\sum_{d\mid n}d\,b_d
\]
be its divisor-ghost sequence.  Define the prime-uniform multiplicative function
\[
 f_b(n)=\prod_{p^e\parallel n}[z^e]A_b(z).
\]
Every prime has the same local factor \(A_b(T)\).  By \cref{thm:block07-euler-classification,thm:block07-product-multiplicative-descent}, the coefficient, exponent, ghost, and prime-uniform multiplicative charts are reversible on this declared locus.

\begin{theorem}[Euler--multiplicative invariant comparison]
\label{thm:block13-em-comparison}
The exponent sequence
\[
 \mathcal J_{\mathrm{EM}}(A_b)=b
\]
is a complete invariant of the certified Euler--multiplicative component.  In every chart it is recovered by the same M\"obius formula
\[
 b_n
 =
 \frac1n
 \sum_{d\mid n}\mu(n/d)c_d,
\]
and Dirichlet convolution corresponds to exponent addition.
\end{theorem}

\begin{proof}
Unique Euler factorization reconstructs \(A_b\) from \(b\).  Prime-power reconstruction then reconstructs \(f_b\).  The logarithmic-derivative identity gives \(c\), and M\"obius inversion recovers \(b\).  Every bridge therefore factors through the same exponent normalizer, so higher coherence follows from \cref{thm:block13-normalizer-coherence}.
\end{proof}

The theorem remains confined to the fixed Euler basis and the prime-uniform multiplicative essential image.  It does not identify arbitrary Cauchy-product grammars with arbitrary Dirichlet-convolution systems.

\subsection{Finite Galois component}
\label{subsec:block13-gal}

Fix a finite Galois extension \(L/K\) in the semisimple setting.  For each subgroup \(H\leq G\), define
\[
 \mathcal J_{\mathrm{Gal}}(H)
 =
 \left(
 H,
 L^H,
 P_H,
 \operatorname{Core}_G(H),
 \operatorname{Ncl}_G(H),
 \{J\backslash G/H\}_{J\leq G}
 \right).
\]

\begin{theorem}[Finite Galois invariant comparison]
\label{thm:block13-gal-comparison}
The subgroup, fixed-field, and faithful Reynolds-projector charts have coherent comparison maps to \(\mathcal J_{\mathrm{Gal}}\).  The invariant is complete on the normalized subgroup atlas for the fixed extension \(L/K\).
\end{theorem}

\begin{proof}
The first three entries determine one another by \cref{thm:block12-normalized-galois-chart}.  Core, normal closure, and double-coset packets are functorial constructions from the subgroup.  The subgroup coordinate itself is complete, and the strict subgroup normalizer forces coherence.
\end{proof}

The Brauer class and Tannakian group scheme are not added to this complete subgroup invariant as if they were automatic coordinates.  They belong to separately typed projective and tensor extensions whose hypotheses and obstructions were stated in Galois-Tannakian atlas.

\subsection{The global tagged invariant}
\label{subsec:block13-global-tagged}

Let
\[
 \Lambda_{13}
 =
 \{\mathrm{DRLM},\mathrm{Loc},\mathrm{Fuss},\mathrm{EM},\mathrm{Gal}\}
\]
be the component-tag set.  Define
\[
 \mathcal J_{13}(X)
 =
 \bigl(\lambda(X),\mathcal J_{\lambda(X)}(X)\bigr).
\]

\begin{theorem}[Certified global completeness theorem]
\label{thm:block13-global-completeness}
The invariant \(\mathcal J_{13}\) descends across every reversible bridge in \(\mathfrak A_{13}\) and is complete on
\[
 \mathcal O_{13}
 =
 \mathcal O_{\mathrm{DRLM}}
 \sqcup
 \mathcal O_{\mathrm{Loc}}
 \sqcup
 \mathcal O_{\mathrm{Fuss}}
 \sqcup
 \mathcal O_{\mathrm{EM}}
 \sqcup
 \mathcal O_{\mathrm{Gal}}.
\]
More precisely,
\[
 \boxed{
 \mathcal J_{13}(X)=\mathcal J_{13}(Y)
 \quad\Longleftrightarrow\quad
 X=Y
 }
\]
inside the declared discrete semantic category.  The associated invariant functor reflects isomorphisms on the descent category.
\end{theorem}

\begin{proof}
Within each component, completeness is \cref{thm:block13-drlm-comparison,thm:block13-local-comparison,cor:block13-fuss-invariant,thm:block13-em-comparison,thm:block13-gal-comparison}.  Their comparison maps are coherent by \cref{thm:block13-normalizer-coherence}.  The component tag prevents an object in one semantic doctrine from being identified with an object in another.  Apply \cref{thm:block13-certified-invariant-descent} componentwise and then take the finite coproduct.
\end{proof}

\begin{remark}[What is genuinely global]
The theorem is not the tautology that every object is classified by itself.  Each component invariant is computed in one native chart and transported through several nonidentical presentation categories.  The global assertion is that these transported coordinates agree coherently and classify the descent object.  The coproduct tag is nevertheless essential: the current Library does not provide reversible bridges joining all five components.
\end{remark}

\subsection{Directional satellites and exact loss records}
\label{subsec:block13-directional-satellites}

The certified reversible core is surrounded by directional bridges that remain useful for theorem transport and computation.  Their failure to enter \(\mathcal J_{13}\) is controlled by explicit loss data.

\begin{table}[htbp]
\centering
\caption{Directional satellites of the certified global invariant atlas.}
\label{tab:block13-directional-satellites}
\small
\begin{tabularx}{\textwidth}{@{}L{0.23\textwidth}L{0.19\textwidth}Y@{}}
\toprule
Directional bridge & Retained structure & Exact loss or obstruction\\
\midrule
Integral regular lattice \(\to\) residue section module & digit transitions and residue semantics & Smith--Tor ghost kernel and missing quotient\\
General regular sequence \(\dashrightarrow\) rational mask & coefficient semantics on the essential image & rational-mask existence and collapse obstruction\\
Algebraic branch \(\to\) holonomic annihilator & D-finite coefficient semantics & selected branch and minimal polynomial may be forgotten\\
Marked nonlinear system \(\to\) semantic fixed point & the represented fixed point & generator labels, arities, and tree grammar may be forgotten\\
Finite-state dynamics \(\to\) Newton--Mahler chart & congruence branch values and precision ladder & depends on a chosen lift and principal-unit hypothesis\\
Projective Galois cocycle \(\dashrightarrow\) linear descent & projective action & relative Brauer class\\
Faithful corner realization \(\dashrightarrow\) Tannakian symmetry & additive packet and operator equations & exact tensor, dual, unit, and fiber-functor data are absent\\
\bottomrule
\end{tabularx}
\end{table}

\begin{proposition}[No completeness transport through an unreflected directional bridge]
\label{prop:block13-no-directional-completeness}
Let \(T:\mathcal C\to\mathcal D\) be a directional bridge.  A complete invariant on \(\mathcal D\) pulls back to a complete invariant on \(\mathcal C\) only if \(T\) reflects isomorphisms on the declared source subcategory.  In particular, any nontrivial object-collapse relation or obstruction fiber prevents automatic completeness transport.
\end{proposition}

\begin{proof}
If \(T\) does not reflect isomorphisms, there exist nonisomorphic \(X,Y\) with \(T(X)\cong T(Y)\).  Every invariant depending only on the target object then takes the same value on \(X\) and \(Y\), so its pullback cannot be complete.
\end{proof}

\subsection{Exact finite verification}
\label{subsec:block13-verification}

Two complementary exact suites audit Global synthesis atlas.  The first authenticates the frozen preceding verification records and checks the broad bridge matrix, periodic and product--multiplicative overlap replays, directional-loss controls, finite \v{C}ech holonomy models, and the five completeness gates.  The second checks the weighted \(q\)-ary tree--algebraic--holonomic overlap, its coefficient and packet equations, parameter recovery, detailed bridge ledger, and tagged component separation.

\begin{table}[htbp]
\centering
\caption{Exact verification of the Global synthesis atlas invariant-descent and weighted q-ary cross-stratum models.}
\label{tab:block13-global-invariant-validation}
\small
\begin{tabularx}{\textwidth}{@{}Y r@{}}
\toprule
Validation class & Checks\\
\midrule
Bridge-matrix structural checks & 172\\
Weighted q-ary coefficient identities & 1,856\\
Exact support-period identities & 8,192\\
Formal fixed-point identities & 8,192\\
Minimal-polynomial relation checks & 8,192\\
Finite parameter recovery & 128\\
Tree/algebraic/holonomic comparison paths & 64\\
Supported first-order recurrence checks & 1,792\\
Packet support and orbit-phase checks & 2,112\\
Exact packet synthesis checks & 44,352\\
Address-carry packet equations & 44,352\\
Tagged component-separation checks & 49\\
\midrule
\textbf{Global synthesis exact checks} & \textbf{119,453}\\
\textbf{Failures} & \textbf{0}\\
\bottomrule
\end{tabularx}
\end{table}


\begin{table}[htbp]
\centering
\caption{Exact verification of the Global synthesis atlas global-invariant and cross-stratum synthesis layer.}
\label{tab:block13-combined-validation}
\small
\begin{tabularx}{\textwidth}{@{}Y r@{}}
\toprule
Validation suite & Exact assertions\\
\midrule
Bridge matrix, overlap replay, directional-loss, and finite \v{C}ech suite & 154,209\\
Weighted $q$-ary tree--algebraic--holonomic and tagged-invariant suite & 119,453\\
\midrule
Regression/authentication assertions already represented in the certified native-stratum program & 39,846\\
\textbf{New Global synthesis atlas comparison and coherence assertions} & \textbf{233,816}\\
All assertions executed by the two Global synthesis atlas suites & 273,662\\
Failures & 0\\
\bottomrule
\end{tabularx}
\end{table}


The two suites execute \(\BlockThirteenChecks\) exact assertions with no failures.  Of these, \(\BlockThirteenRegressionChecks\) authenticate or replay assertions already represented in the preceding cumulative record.  The global synthesis therefore contributes \(\BlockThirteenNewChecks\) nonduplicative comparison, coherence, packet, and loss-control assertions.  The blockwise cumulative record is
\[
 \boxed{\BlockThirteenCumulative}
\]
exact assertions with no recorded failures.  Counting every regression assertion executed during the Global synthesis atlas replay gives \(\BlockThirteenCumulativeExecuted\) total executed assertions.  Neither number is a count of independent theorems.

The bridge-synthesis suite replays primitive unit-periodic words over finite prime fields, compares their recurrence, OGF, automaton, section-module, and functional-Mahler coordinates, verifies the prime-uniform Euler--multiplicative overlap, and tests finite nontrivial \v{C}ech cycle models.  The weighted \(q\)-ary suite checks the fixed-point identity, Lagrange coefficients, first-order P-recursive recurrence, support phases, exact packet synthesis, address-carry equations, and all triangle normalizers over a declared finite parameter grid.  Both machine-readable records are included in the source package.

\begin{boundarybox}
\textbf{Exact boundary of Global synthesis atlas.}
The global invariant theorem is proved for the five-component certified finite atlas in \cref{def:block13-certified-atlas}.  It does not assert a complete invariant for every object in the eight broad strata, a reversible bridge between the five semantic components, effective descent for every directional satellite, cross-radix or all-ring normalization, modular or wild Galois effectivity, nonneutral Tannakian reconstruction, or classification of unmarked general holonomic, algebraic, and nonlinear systems.
\end{boundarybox}

\begin{resultbox}
\textbf{Section result.}
The manuscript now contains an authoritative broad-stratum bridge matrix, a finite certified atlas with five nontrivial connected components, explicit invariant comparison maps, a normalizer-induced higher-coherence theorem, and a complete tagged global invariant on the declared semantic union.  The new Fuss--Catalan triangle gives a genuine nonlinear--algebraic--holonomic overlap.  The digital--regular--functional-Mahler, continuous Newton--Mahler, Euler--multiplicative, and finite Galois components consolidate the exact cross-presentation results of the preceding stratum sections without conflating the two Mahler theories.  Directional bridges remain outside the reversible core and retain their precise loss or obstruction records.  This closes the finite global-invariant milestone on a substantial atlas, not the unrestricted eight-stratum stack.
\end{resultbox}
